\section{Exercises}\label{sec:12-path-algebras:06-exercises:1}

\textbf{Key points.} A quiver is just vertices, arrows, and source/target functions; a path is an inductive type indexed by its own endpoints, so an ill-typed composition (arrow source ≠ current endpoint of the path) is rejected before any proof is attempted. \texttt{Path.append} composes two paths by recursion on the second, and the path algebra \(kQ\), sketched but not fully built here, is the free \(k\)-module on the set of paths, with multiplication by composition.

\textbf{Socratic questions.}

\begin{enumerate}
\def\labelenumi{\arabic{enumi}.}
\item
  \emph{The composability check of \texttt{Quiver.Path} (the source of an arrow must match the current endpoint of the path) is enforced by Lean at compile time. What would the equivalent check look like in an ordinary, non-dependently typed program, and when would it fire?} At runtime, as an \texttt{if}-check or assertion inside whatever function tries to append, discovered the moment that specific composition executes, not before. Here, an ill-typed \texttt{Path.cons} simply has no proof term to supply, so the program that would have contained the bug never compiles at all.
\item
  \emph{A quiver with a genuine cycle (the \texttt{gamma} of Exercise 1) has infinitely many distinct paths (go around the loop any number of times). Does that mean \texttt{Path Q u v} is an infinite }type\emph{, for such a \texttt{Q}?} For at least one pair \texttt{u = v} on the cycle, yes. Infinitely many distinct terms inhabit that one type, one per lap. This is exactly why the path algebra \(kQ\) of a quiver is typically infinite-dimensional once the quiver has a cycle, unless further relations are imposed.
\item
  \emph{The path algebra \(kQ\) is described as ``the free \(k\)-module on the set of paths,'' but this chapter builds only the paths and their composition, not \(kQ\) itself. What is the one piece of extra machinery still missing?} Finitely-supported functions from paths to \(k\), the ability to form a \(k\)-linear combination of \emph{finitely many} paths at once, which is genuinely more than what an inductive \texttt{Path} type alone provides.
\item
  Add a third arrow \texttt{gamma : ExampleArrow} with \texttt{source gamma = 2} and \texttt{target gamma = 0}, creating a cycle \texttt{0 → 1 → 2 → 0}. Build the path \texttt{gamma ∘ beta ∘ alpha : Path exampleQuiver 0 0}.
\item
  Prove \texttt{theorem append\_\allowbreak{}nil\_\allowbreak{}left \{V A\} \{Q : Quiver V A\} \{u v\} (p : Path Q u v) : Path.append (Path.nil u) p = p} by induction on \texttt{p} (mirroring the structure of the recursion of \texttt{Path.append} itself), spelling out the base (\texttt{Path.nil}) and inductive (\texttt{Path.cons}) cases as in the \texttt{induction} examples of Chapter 5.
\item
  (Optional, harder) Sketch, in comments, no need to complete the Lean, what a \texttt{structure PathAlgebra (V A : Type) (Q : Quiver V A) (k : Type) (Rg : Ring k)} would need to contain to satisfy the fields of \texttt{Ring} from Chapter 9.
\end{enumerate}

Solutions, \hyperref[sec:15-appendix-solutions:12-chapter-12:1]{Appendix, Chapter 12}.

